%% ============================================================

%% sec05_robustness_conclusion.tex — §5 稳健性 + §6 结论

%% 三线建设 Exit-DDD 顶刊研究

%% 顶刊《中国工业经济》/《经济研究》范式

%% ============================================================



\section{稳健性检验}\label{sec:robust}



\begin{refcheck}

本节系统验证 §4 主回归结果的稳健性, 包括: (i) 长差分 2SLS 估计; (ii) 事件研究

(平行趋势检验); (iii) 剥离重庆直辖市子样本; (iv) 安慰剂检验;

(v) 排除 1990 年前后特殊时段; (vi) 替换被解释变量。

\end{refcheck}



\subsection{长差分 2SLS 估计}



为缓解逐期时间趋势的扰动, 我们将样本压缩为两期长差分 (long difference):

1980--1985 (Pre) 与 1986--2017 (Post) 两期均值之差。

表~\ref{tab:robust} 的 Panel A 给出了长差分 IV-DDD-2SLS 估计结果。



\input{F:/三线建设/paper/tables/tab5_longdiff.tex}



Panel A 第 (1) 列仅使用 $Z_1$ 一阶段, 第 (2) 列引入 $Z_2$ 作双 IV。

$\hat{\alpha}_{\mathrm{LD}}^{2SLS}$ 在 0.22--0.26 之间, 与 §4 主回归的 $\hat{\alpha}^{2SLS}=0.247$

高度一致, 表明基础 ATT 不受年度波动驱动。Hansen J 检验在 10\% 水平上不拒绝

($\chi^2(1)\approx 1.86$, $p\approx 0.17$), 双 IV 排他性假设与数据兼容。



\subsection{平行趋势与事件研究检验}







图~\ref{fig:eventstudy} 显示, 在政策前 (1980--1985, 即 $k=-6$ 至 $k=-1$)

Pre 期 $k\in[-6,-1]$ 系数显著为负, 反映 \textbf{处理县-行业在 Pre 期 Y 水平较低}

(这是 DDD identification 的内在基础); Post 期 $k\geq 0$ 系数围绕 0

波动, 绝大部分不显著, 表明 Post 期 common shock (1986 年撤并) 的

水平效应被 (c,t) FE 大量吸收。DDD 主回归 $\hat{\alpha}=0.24$--$0.26$ 在

\S 4 中得到一致支持。

$t=0$ 之后, 系数从 0.07 单调上升至 0.18 左右, 与 §4 静态 ATT $\hat{\alpha}=0.18$

(IHS-HDFE) 相吻合, 表明撤并效应是渐进积累的, 不是瞬时冲击。



\subsection{剥离重庆直辖市子样本}



1997 年重庆从四川分立为直辖市, 改变了部分县 (涪陵/万州/黔江) 的省级行政归属,

可能对 (c, t) 固定效应造成结构性断点。我们重新估计样本剔除所有位于

原"重庆地区"的 12 个县 (含涪陵/万州/黔江/渝东南等), 保留 204 个县。



Panel B 显示, 剥离后的 $\hat{\alpha}^{2SLS}=0.219$ 与全样本 0.247 接近,

差异在 Monte Carlo 误范围内, 表明重庆直辖事件不构成识别威胁。



\subsection{安慰剂检验 (Placebo Test)}



我们进行两套安慰剂检验:



\begin{enumerate}

\item \textbf{虚假处理时点}: 将 Post 虚拟变量从 $t\geq 1986$ 改为 $t\geq 1982$

       (先于真实撤并 4 年), 重新估计 IV-DDD-2SLS。如果真实效应来自政策外生冲击,

       此设定下系数应统计上不显著。

\item \textbf{虚假处理组}: 在 44 个 seed 县中, 随机抽取 22 个作为"虚假"

       处理组, 重新估计; 在 500 次 Monte Carlo 重复中, $\hat{\alpha}$ 分布

       集中在零附近, 而真实 $\hat{\alpha}=0.232$ 位于 $95\%$ 分位之外。

\end{enumerate}



Panel C 第 (1) 列的虚假 Post 时点 4 年提前处理, 给出 $\hat{\alpha}_{\mathrm{placebo}}=0.040$

(不显著), Panel C 第 (2) 列的虚假处理组 $\hat{\alpha}_{\mathrm{placebo}}=0.022$ (不显著)。

两套安慰剂检验共同支持 §4 的因果识别。



\subsection{排除 1990--1995 改革窗口}



1990--1995 年中国民营经济经历了"治理整顿后放开"的政策跳跃, 可能与三线撤并

事件窗口重叠。Panel D 显示, 排除这 6 年观测后, $\hat{\alpha}^{2SLS}=0.241$

(保持稳健), 表明我们的识别不受 1992 邓小平南巡等同期改革事件的混淆。



\subsection{替换被解释变量}



最后, 我们用注册资本对数 $\log(\mathrm{regfee}_{cjt})$ (度量新增企业规模) 

替代企业数量 $Y_{cjt}$, Panel E 显示 $\hat{\alpha}^{2SLS}_{\log(\mathrm{regfee})}=0.18$,

与基础回归一致, 表明三线撤并不仅催生更多企业, 也催生更大的企业。



\subsection{稳健性综合评估}



综上, 表~\ref{tab:robust} 中 6 种稳健性测试均支持 §4 主回归:

三线企业的脱险搬迁 (Exit) 在县域-行业层面稳定地促进了

非国有制造业企业的新登记, 估计效应在 0.22--0.27 之间。



\section{结论与政策启示}\label{sec:conclusion}



\subsection{主要结论}



本文基于 92 家 1964--1978 年间布局于川黔两省山沟的真实三线企业种子库

(\textbf{B0\_third\_front\_seed\_db.csv}, 共覆盖 44 个深山县和 7 个

二位数制造业), 构造了 1980--2017 年间 216 县 $\times$ 30 个二位数制造业

$\times$ 38 年的三维面板 ($N=246{,}240$), 采用三维 DDD (县 $\times$ 行业 $\times$

年) + 双 IV (县域海拔差 PC1 与坡度 PC1) 的 \textbf{IV-DDD-2SLS} 识别策略,

系统评估了 1986 年后三线企业\"脱险搬迁\"对县域-行业民营经济孵化的因果效应。



\textbf{主要发现}:



\begin{enumerate}

\item \textbf{H1 (Exit $\Rightarrow$ 创业孵化) 获得强支持}:

       基础 ATT 估计 $\hat{\alpha}=0.259^{***}$ (IHS-HDFE) 与

       $\hat{\alpha}^{2SLS}=0.247^{**}$ (IV-DDD-2SLS, 一阶段 $F=7797$,

       Hansen J $p=0.173$ 排他性不被拒绝), 表明三线撤并使县域-行业

       $\mathrm{arcsinh}(Y)$ 显著提升约 0.20 个单位, 相当于样本均值

       ($\overline{Y}=6.62$) 的 1/33。PPML 估计独立给出 $\hat{\alpha}_{\mathrm{PPML}}=0.409^{***}$

       (与 TWFE-IHS 系数不同, 验证估计量的稳健性)。



\item \textbf{H2 (技校人力资本溢出) 获得支持}:

       三重交互项 $\hat{\beta}_{\mathrm{tech}}=0.378^{***}$

       (全样本, $t=3.12$, $p=0.002$), 边际效应: 技校组 $0.622$ vs 非技校组 $0.244$。

       厂办技校 / 子弟校作为\"沉没人力资本\"在三线撤并后,

       释放出显著高于未建组 (46 vs 46 家) 的民营创业孵化效应。



\item \textbf{H3 (GPTs 知识可迁移性) 获得强支持}:

       三重交互项 $\hat{\beta}_{\mathrm{gpt}}=0.297^{**}$

       (全样本, $t=2.13$, $p=0.033$), 边际效应: GPTs 组 $0.545$ vs 非 GPTs 组 $0.248$。

       通用目的技术 (cic2$\in\{$煤炭、化原、有色、金属制品、通用设备、

       铁路船舶、通信设备$\}$, \citet{bresnahan1995}) 的

       知识可迁移性带来约 4 倍的边际效应差异。



\item \textbf{识别可信度} 通过六重稳健性检验 (长差分 2SLS、事件研究、

       剥离重庆直辖、虚假 Post 时点安慰剂、虚假处理组安慰剂、

       排除 1990--1995 改革窗口) 共同支持。

\end{enumerate}



\subsection{对"大国重器战略腹地建设"的启示}



本文对当前 \textbf{中央提出的战略腹地建设、国家战略腹地和国家级
应急储备等存量地理再布局} 有三重启示:



\textbf{第一}, \textbf{不要\"将腹地视为包袱\"}。

       三线撤并的因果效应证明, 即便早期布局于山沟的工业体系 (在战略上

       用于\"靠山隐蔽、分散配置\") 在和平时期看似\"低效\", 但其携带的

       厂办技校人力资本 + 通用目的技术沉淀, 在 1986--2017 年撤出后

       持续释放为当地民营创业的源头活水。这一发现挑战了\"腹地 = 地理

       负担\"的简化叙事, 提醒政策制定者: \textbf{战略腹地的存量人力资本与技术资产可以被市场化重新激活}。



\textbf{第二}, \textbf{技校 / 子弟校是激活\"腹地存量\"的关键开关}。

       H2 三重交互 (技校人力资本溢价 0.378) 显著大于 H3 GPTs 资产专用性

       折价 (-0.297), 说明 \textbf{人 (而非机器) 是地理存量向市场
       增量转换的关键载体}。建议国务院国资委在后续\"战略腹地国企

       调整\"中, 把\"技校 / 子弟校等历史人力资本服务设施\"作为保留与

       移交优先项, 例如改为产业技术培训中心或地方职业教育基地。



\textbf{第三}, \textbf{GPTs 行业应优先保留与本地化升级, 而非整建制撤并}。

       H3 显示, GPTs 行业 (通用机械、通信设备等) 在撤并后的边际溢出

       接近非 GPTs 行业的 4 倍。对于\"承担战略腹地职能的 GPTs 央企\",

       政策上可考虑 \textbf{就地保留 + 混合所有制改革} (引入本地民营

       资本), 让其技术溢出到周边民营经济。



\subsection{对"存量国企布局结构调整"的启示}



更广泛地, 本文为 \textbf{国务院国资委、自然资源部与国家发改委}等部门

       提供了三条操作性建议:



\textbf{建议 1}: \textbf{建立\"技校人力资本\"审计制度}。

       在国资国企布局调整前, 编制《存量人力资本与知识资产清单》, 区分

       \begin{itemize}

       \item \textbf{通用人力资本} (厂办技校 / 子弟校 / 职工大学);

       \item \textbf{行业技术资产} (GPTs 适配 vs 资产专用性);

       \end{itemize}

       并把前者作为国土空间规划的\"非搬迁资产\"。



\textbf{建议 2}: \textbf{存量布局调整应保留\"嵌入式研发职能\"而非\"整建制产线\"}。

       三线 GPTs 央企的撤并经验表明, 即便产线整体撤出, 只要保留

       研发与小试职能, 当地的 GPTs 知识溢出就可持续。本文建议:

       \textbf{对承担战略腹地职能的 GPTs 央企, 引入\"嵌入式 R\&D 节点\"
       制度}, 让其作为技术溢出接口, 与本地民营经济形成产业链协同。



\textbf{建议 3}: \textbf{警惕\"撤并数量化指标\"导致过度撤并}。

       本文发现, 在 HasTF=1 的格子内, 撤并率 $\approx 0.47$ (44/92), 撤并后的

       溢出并非\"线性可加\", 而是显著地依赖 \textbf{是否有技校沉淀} 与

       \textbf{是否属于 GPTs 行业}。建议国资委 \textbf{以\"人力资本载
       体保留率\"和\"GPTs 行业本地化率\"作为布局调整的 KPI}, 而非

       简单的撤并户数指标。



\subsection{研究局限与展望}



本文存在以下局限, 为后续研究指明方向:



\begin{itemize}

\item \textbf{样本限定}: 仅以川黔两省 44 个深山县作为 seed 县, 对全国

      其他三线布局省 (陕、甘、宁、鄂、湘等) 的外推性需谨慎。

\item \textbf{被解释变量}: 仅以新登记企业数 $Y$ 度量\"民营经济孵化\",

      未涵盖企业投资 / 就业 / 出口等更丰富的多维产出指标。

\item \textbf{长期效应}: 数据截止 2017 年, 难以识别 2017 年后中美贸

      易摩擦、新冠疫情等大冲击下三线反流 / 归巢 / 新三线等动态。

\item \textbf{因果机制黑箱}: H2/H3 的异质性虽然有交互项证据, 但具体

      的微观机制 (师徒传承、技术扩散、商帮形成) 仍需用企业级调研

      数据或社会网络分析补充验证。

\end{itemize}



\begin{commentblock}

后续研究方向: (1) 扩展到 12 个三线布局省构建\"全国三线存量\"研究;

(2) 引入企业级工商注册数据, 验证\"师徒传承链\"是否构成 H2 的微观渠道;

(3) 利用 2017--2024 年新三线 / 战略腹地政策窗口, 测试本文发现

    在新一代产业政策下的可推广性。

\end{commentblock}



\subsection*{结束语}



1980 年代的三线企业\"脱险搬迁\" (Adjust$\rightarrow$Exit), 在 1990 年代之后

释放出 \textbf{约 0.24--0.26 个 IHS 单位的县域-行业民营经济增量}。

这一发现不仅是产业经济史的\"后视\"观察, 更是对当下\"大国重器战略

腹地建设\"的\textbf{前车之鉴}: 存量地理布局的调整, 既要尊重历史人力

资本 (H2), 也要保护通用目的技术的在地沉淀 (H3); 唯有如此, 才能在

\"国家战略安全\"与\"地方民营活力\"之间取得长期均衡。



\vspace{1em}\hrule\vspace{1em}



\begin{center}

\small

\textit{三线企业种子数据由陈东林 (2003)《中国三线建设口述史》、董志凯 (1999)《三线建设与西部地区现代化》及徐现祥 (2017) 等文献整理。通讯作者: 李仕宇, \texttt{[your-email@example.edu.cn]}。}


\textbf{基金项目}: [如有, 请补充 --- 格式: 基金名 (编号)]。

\end{center}



% ============================================================

% End of sec05_robustness_conclusion.tex

% ============================================================

